% !TeX program = xelatex
\usepackage[b5paper,left=22mm,right=22mm,top=20mm,bottom=22mm,footskip=11mm]{geometry}
\usepackage{fontspec}
\setmainfont{TeX Gyre Termes}
\setsansfont{TeX Gyre Heros}
\usepackage{amsmath,amssymb,mathtools,amsthm}
\renewcommand{\cdots}{\mathinner{\cdotp\cdotp\cdotp}}
\usepackage{graphicx,adjustbox,array,booktabs,longtable,multirow,enumitem,changepage,needspace}
\usepackage[table]{xcolor}
\usepackage{fancyhdr}
\raggedbottom
\interfootnotelinepenalty=10000
\newcommand{\BookFootnote}[2]{\begingroup\renewcommand{\thefootnote}{\fnsymbol{footnote}}\footnote[#1]{#2}\endgroup}
\newcommand{\BookFootnoteMark}[1]{\begingroup\renewcommand{\thefootnote}{\fnsymbol{footnote}}\footnotemark[#1]\endgroup}
\newcommand{\BookFootnoteText}[2]{\begingroup\renewcommand{\thefootnote}{\fnsymbol{footnote}}\footnotetext[#1]{#2}\endgroup}
\newenvironment{BookFloat}[1]{\par\addvspace{.6\baselineskip}\noindent\def\BookFloatID{#1}\begin{minipage}{\linewidth}\phantomsection\label{bookfloat:#1:start}}{\phantomsection\label{bookfloat:\BookFloatID:end}\end{minipage}\par\addvspace{.6\baselineskip}}

\usepackage[unicode,hidelinks]{hyperref}
\pagestyle{fancy}
\fancyhf{}
\newcommand{\BookFooter}{}
\fancyfoot[L]{\parbox[b]{\dimexpr\textwidth-15mm}{\footnotesize\BookFooter}}
\fancyfoot[R]{\footnotesize\thepage}
\renewcommand{\headrulewidth}{0pt}
\renewcommand{\footrulewidth}{0pt}
\setlength{\parindent}{0pt}
\setlength{\parskip}{0.2em}
\setlength{\emergencystretch}{3em}
\setlist{nosep,leftmargin=1.7em,labelsep=.65em}
\newif\ifBookNarrative
\BookNarrativetrue
\newenvironment{BookBody}{\setlength{\parindent}{0pt}}{}
\newcommand{\BookStatement}[1]{\par\addvspace{.6\baselineskip}\noindent #1\enspace\ignorespaces}
\newcommand{\BookLabel}[1]{\par\addvspace{.6\baselineskip}\needspace{3\baselineskip}\noindent\textbf{#1}\quad\ignorespaces}
\newcommand{\BookHeading}[1]{\par\addvspace{\baselineskip}\needspace{4\baselineskip}\noindent{\large\bfseries #1}\par\medskip}
\newcommand{\BookSubheading}[1]{\par\addvspace{.7\baselineskip}\needspace{3\baselineskip}\noindent{\bfseries #1}\par\smallskip}
\newcommand{\BookTextHeading}[1]{\par\addvspace{.7\baselineskip}\needspace{3\baselineskip}\noindent #1\par\smallskip}
\newcommand{\BookSectionTitle}[2]{\par\addvspace{1.2\baselineskip}\needspace{5\baselineskip}\noindent\colorbox{gray!25}{\sffamily\bfseries\small SECTION #1}\quad{\large\bfseries #2}\par\medskip}
\newcommand{\BookExerciseTitle}[1]{\par\addvspace{\baselineskip}\needspace{4\baselineskip}\noindent{\color{gray}\rule{1.5mm}{1.5mm}}\enspace{\large\bfseries Exercises #1}\par\medskip}
\newcommand{\BookBullet}[1]{\begin{list}{\textbullet}{\setlength{\leftmargin}{1.7em}\setlength{\labelwidth}{1em}\setlength{\labelsep}{.7em}\setlength{\topsep}{0pt}\setlength{\parsep}{0pt}\setlength{\itemsep}{0pt}}\item #1\end{list}}
\newcommand{\BookExercise}[2]{\begin{list}{\textbf{#1}}{\setlength{\leftmargin}{2.2em}\setlength{\labelwidth}{1.5em}\setlength{\labelsep}{.7em}\setlength{\topsep}{.15em}\setlength{\parsep}{0pt}\setlength{\itemsep}{0pt}}\item #2\end{list}}
\newcommand{\BookEnd}[1]{\unskip\nobreak\hfill\ensuremath{#1}}
\newcommand{\BookSource}[2]{\phantomsection\label{sourcepage:#1}} % Source-page anchors, retained through reflow.
\usepackage{multicol}
\usepackage{tikz}
\newcommand{\BookPartEntry}[2]{\par\noindent\makebox[17mm][l]{\bfseries Section #1}#2\par}
\newcommand{\BookPartTitle}[3]{\par\bigskip\noindent\begin{tikzpicture}\node[anchor=north west,inner sep=0pt,outer sep=0pt] (bookpart) at (0,0) {\begin{minipage}{\dimexpr\linewidth-.5pt}\hspace*{26mm}\begin{minipage}[t]{\dimexpr\linewidth-26mm}\begin{adjustbox}{max width=\linewidth}{\fontsize{20}{24}\selectfont\bfseries #2}\end{adjustbox}\par\vspace{6mm}\hrule\vspace{6mm}{\sffamily\footnotesize\setlength{\parskip}{.3em}#3}\vspace{3mm}\end{minipage}\end{minipage}};\path[shade,top color=gray!60,bottom color=white] (0,0) rectangle ([xshift=24mm]bookpart.south west);\node[anchor=north,font=\sffamily\bfseries\small] at (12mm,-2mm) {PART};\node[anchor=north,font=\sffamily\bfseries\Huge] at (12mm,-7mm) {#1};\end{tikzpicture}\par\bigskip}
\newsavebox{\BookMathBox}
\makeatletter
\newcommand{\BookInline}[1]{\leavevmode\hskip0pt\relax\begingroup\in@{\begin}{#1}\ifin@\sbox{\BookMathBox}{\(#1\)}\ifdim\wd\BookMathBox>\dimexpr\linewidth-1em\relax\resizebox{\dimexpr\linewidth-1em\relax}{!}{\usebox{\BookMathBox}}\else\(#1\)\fi\else\in@{\left}{#1}\ifin@\sbox{\BookMathBox}{\(#1\)}\ifdim\wd\BookMathBox>\dimexpr\linewidth-1em\relax\resizebox{\dimexpr\linewidth-1em\relax}{!}{\usebox{\BookMathBox}}\else\(#1\)\fi\else\(#1\)\fi\fi\endgroup\hskip0pt\relax}
\makeatother
\DeclareMathOperator{\Aut}{Aut}
\DeclareMathOperator{\End}{End}
\DeclareMathOperator{\Hom}{Hom}
\DeclareMathOperator{\GF}{GF}
\DeclareMathOperator{\irr}{irr}
\DeclareMathOperator{\lt}{lt}
\DeclareMathOperator{\lp}{lp}
\providecommand{\textcircled}[1]{\textnormal{(#1)}}
